
\subsubsection{Multi-Dimensional Trustworthiness Evaluation}

Recent frameworks have advanced multi-dimensional trustworthiness evaluation infrastructure. TrustVis \citep{Sun2025} introduces a unified framework measuring safety and robustness together using adversarial perturbations. MLLMGuard \citep{Gu2024} reports multi-dimensional safety scores across five dimensions (Privacy, Bias, Toxicity, Truthfulness, Legality) using the GuardRank framework. BOLD \citep{dhamala2021bold} benchmarks fairness metrics on 23,679 prompts spanning demographic dimensions. These frameworks generate evaluation logs that are inherently multi-dimensional but stop at per-dimension score reporting without computing or testing cross-dimensional correlations. The present work extends this foundation by analyzing correlation structure.

TrustVis performs sequential assessment with adversarial modification between dimensions, making correlations reflect perturbation effects rather than inherent dimensional coupling. The present methodology differs by measuring dimensions on the same natural inputs without perturbation, isolating genuine correlations from evaluation artifacts. This represents an analytical extension: where TrustVis reports separate safety and robustness scores, the present work computes Pearson correlations and tests for statistical significance.

\subsubsection{Independent Dimension Benchmarks}

Traditional trustworthiness benchmarks evolved in parallel communities. TruthfulQA \citep{lin2022truthfulqa} tests factual accuracy on 817 questions designed to elicit common misconceptions. AdvGLUE \citep{wang2021adversarial} and adversarial suffix attacks \citep{zou2023universal} measure consistency under perturbation. HONEST \citep{nozza2021honest} quantifies demographic bias in language model continuations. The present study uses TruthfulQA as the base dataset, extending evaluation to include robustness and fairness on the same outputs.

\subsubsection{Alignment and Safety Trade-offs}

The term ''alignment tax'' appears informally in RLHF literature to describe safety interventions that reduce model capabilities. Bai et al. (2022) observe that Constitutional AI training can decrease performance on benign tasks. Ouyang et al. (2022) note that instruction-following fine-tuning sometimes reduces factual accuracy. Prior work reports these trade-offs anecdotally or through separate ablations rather than as correlation magnitudes. The present study provides a quantitative estimate (r=-0.25) of the fairness-reliability correlation.

\subsubsection{Training Dynamics and Memorization}

Carlini et al. (2021) demonstrate that large language models memorize and extract verbatim training data. Brown et al. (2020) observe that scaling model size increases memorization capacity, particularly for factual knowledge. The present work extends this by showing that reliability and robustness correlate on factual content (r=0.72), interpreted as consistent with shared memorization mechanisms. The mechanism specificity---factual prompts show r=0.72 versus misinformation prompts showing weaker correlation---provides convergent evidence for this interpretation.
